\begin{lemma}
\label{lem:p071-terminal-polyhedral-insertion}
Let \(Q\) be a nonempty compact convex polytope, let \(R\in\mathbf R\), and
put
\[
                       L=[-\infty,R]\times Q.
\]
Suppose that \(u:L\to\mathbf R\cup\{-\infty\}\) is upper semicontinuous,
jointly extended-real convex on the finite stratum, nondecreasing in the
first variable, and has the exact pointwise trace at \(-\infty\).  If
\(\phi\in C(L)\) and \(u\leq\phi\), then there is a continuous real-valued
function \(E_\phi\) on \(L\), jointly convex on the finite stratum,
nondecreasing in the first variable, and having the exact endpoint trace,
such that
\begin{equation}
                             u\leq E_\phi\leq\phi.
 \label{eq:p071-terminal-insertion-sandwich}
\end{equation}
Moreover, \(E_\phi\) is the largest function with these properties below
\(\phi\), and the operator \(\phi\mapsto E_\phi\) is nondecreasing and
nonexpanding in the uniform norm.
\end{lemma}

\begin{proof}
For a finite point \((t,y)\), let
\begin{equation}
 E_\phi(t,y)=
 \inf\left\{
   \sum_{i=1}^q\alpha_i\phi(t_i,y_i):
   \begin{array}{l}
    q<\infty,\quad \alpha_i\geq0,\quad \sum_i\alpha_i=1,\\
    t_i\in\mathbf R,\quad t_i\leq R,\quad y_i\in Q,\\
    \sum_i\alpha_i t_i\geq t,\quad
    \sum_i\alpha_i y_i=y
   \end{array}\right\}.
 \label{eq:p071-terminal-monotone-convex-envelope}
\end{equation}
At the endpoint, put
\begin{align}
 m(y)&=\min_{s\in[-\infty,R]}\phi(s,y),\\
 E_\phi(-\infty,y)
   &=\inf\left\{\sum_{i=1}^q\alpha_i m(y_i):
       q<\infty,\ \alpha_i\geq0,\ \sum_i\alpha_i=1,\
       \sum_i\alpha_i y_i=y\right\}.
 \label{eq:p071-terminal-endpoint-envelope}
\end{align}
The second expression is the ordinary convex envelope of \(m\) on \(Q\).
The function \(m\) is continuous, since uniform continuity of \(\phi\) on
\(L\) controls its minimum uniformly in the parameter \(y\).

The singleton measure shows that \(E_\phi\leq\phi\) at every point.  Mixtures
of admissible families in
\eqref{eq:p071-terminal-monotone-convex-envelope} show that \(E_\phi\) is
jointly convex on the finite stratum.  Enlarging the lower bound imposed on
the mean of the \(t_i\) shrinks the admissible family, so \(E_\phi\) is
nondecreasing in \(t\).

We next prove continuity away from the endpoint.  For \(N>-R\), restrict
\eqref{eq:p071-terminal-monotone-convex-envelope} to
\(t_i\in[-N,R]\), and denote the resulting envelope by \(E_{\phi,N}\).
For this truncated problem the same infimum is obtained if arbitrary Borel
probability measures on \([-N,R]\times Q\) are allowed.  Indeed,
Carath\'eodory's theorem, applied to the vector consisting of the
\(t\)-coordinate, the affine coordinates on \(Q\), and the value of \(\phi\),
replaces any such measure by a finitely supported measure with the same two
moments and the same cost.  The full probability-measure space is compact in
the weak topology, and the moment conditions are closed.  Hence a convergent
sequence of roots and minimizing measures proves lower semicontinuity.  The
function is finite and convex on the compact polytope \([-N,R]\times Q\), so
a triangulation having any prescribed boundary point as a vertex gives upper
semicontinuity there.  Thus \(E_{\phi,N}\) is continuous.

Set
\[
 \omega_N=\sup\{
   |\phi(s,y)-\phi(s',y)|:
   s,s'\in[-\infty,-N],\ y\in Q\}.
\]
Continuity on the compactification gives \(\omega_N\to0\).  Since the
restricted admissible family is smaller,
\[
                         E_\phi\leq E_{\phi,N}
                         \quad\text{on }[-N,R]\times Q.
\]
Conversely, in any admissible finite family for \(E_\phi(t,y)\), replace
each \(t_i<-N\) by \(-N\).  This only increases its \(t\)-mean, preserves
its \(Q\)-barycentre, and changes its cost by at most \(\omega_N\).
Taking infima gives
\begin{equation}
 0\leq E_{\phi,N}-E_\phi\leq\omega_N
 \quad\text{on }[-N,R]\times Q.
 \label{eq:p071-terminal-finite-truncation}
\end{equation}
On every fixed finite subprism, the right side tends to zero.  Hence
\(E_\phi\) is continuous at every finite point.

The endpoint formula is also continuous in \(y\): the convex envelope of a
continuous function on a compact polytope is finite and continuous.  Every
family admissible in
\eqref{eq:p071-terminal-monotone-convex-envelope} gives
\[
 E_\phi(-\infty,y)
 \leq\sum_i\alpha_i m(y_i)
 \leq\sum_i\alpha_i\phi(t_i,y_i),
\]
and therefore \(E_\phi(-\infty,y)\leq E_\phi(t,y)\).
For the reverse estimate, choose a minimizing finite representation of
\(E_\phi(-\infty,y)\).  It exists by the same compactness and Carath\'eodory
argument, now on \(Q\).  For each \(y_i\), choose a point \(s_i\) where
\(m(y_i)=\phi(s_i,y_i)\).  Replace every
\(s_i<-N\), including \(s_i=-\infty\), by \(-N\).  The resulting family
has \(t\)-mean at least \(-N\), and its cost increases by at most
\(\omega_N\).  Consequently
\begin{equation}
 0\leq E_\phi(t,y)-E_\phi(-\infty,y)\leq\omega_N
 \quad(t\leq-N,\ y\in Q).
 \label{eq:p071-terminal-uniform-endpoint}
\end{equation}
This proves both the exact trace and joint continuity at the endpoint.

It remains to prove the first inequality in
\eqref{eq:p071-terminal-insertion-sandwich}.  For a finite root and a family
in
\eqref{eq:p071-terminal-monotone-convex-envelope}, monotonicity followed by
convexity gives
\[
 u(t,y)\leq
 u\left(\sum_i\alpha_i t_i,\sum_i\alpha_i y_i\right)
 \leq\sum_i\alpha_i u(t_i,y_i)
 \leq\sum_i\alpha_i\phi(t_i,y_i).
\]
By the exact trace, the endpoint slice is the pointwise decreasing limit, as
\(t\to-\infty\), of the convex slices \(u(t,\cdot)\); it is therefore
extended-real convex.  At the endpoint, monotonicity and this convexity give
\[
 u(-\infty,y)\leq
 \sum_i\alpha_i u(-\infty,y_i)
 \leq\sum_i\alpha_i m(y_i).
\]
Taking the two infima proves \(u\leq E_\phi\).

The same two Jensen calculations apply to every jointly convex,
nondecreasing function with exact trace below \(\phi\).  Such a function is
therefore bounded above by \(E_\phi\).  Since the preceding paragraphs show
that \(E_\phi\) itself has all the required properties, it is the largest
one.  Monotonicity in the obstacle is immediate from the two infimum
formulas.  If \(\lVert\phi-\psi\rVert_\infty\leq\delta\), every admissible
cost for one obstacle differs from the corresponding cost for the other by
at most \(\delta\).  Hence
\(\lVert E_\phi-E_\psi\rVert_\infty\leq\delta\).
\end{proof}
